\chapter{VAR Estrutural}
\label{cap:svar}
% ── índice ──────────────────────────────────────────────────────
\index{VAR estrutural|textbf}
\index{SVAR|see{VAR estrutural}}
\index{svar@\texttt{svar()}|textbf}
\index{sirf@\texttt{sirf()}|textbf}
\index{Cholesky, decomposição de|textbf}
\index{identificação estrutural|textbf}
\index{função impulso-resposta estrutural|textbf}
\index{ordenação causal}

% ─────────────────────────────────────────────────────────────────────────────
\section{Da forma reduzida à estrutura}

O VAR da forma reduzida (cap.~\ref{cap:var}) fornece $K^2 \cdot p + K$ parâmetros
livres. A impulso-resposta reduzida mede o efeito de \emph{inovações} $u_t$
sem conteúdo econômico. O SVAR recupera \emph{choques estruturais}
$\varepsilon_t$ (oferta, demanda, política monetária) impondo restrições
que refletem teoria econômica:

\[
  A\,y_t = A_1^*\,y_{t-1} + \cdots + B\,\varepsilon_t
  \qquad
  \varepsilon_t \sim (0, I_K)
\]

\subsection*{Identificação de Cholesky}

A decomposição triangular inferior de $\Sigma_u = P\,P'$ impõe uma
ordenação causal recursiva: a variável na posição 1 é exógena
contemporaneamente; a na posição 2 responde contemporaneamente apenas
à primeira; etc.

\[
  B = P \quad\text{(triangular inferior)} \implies y_{1,t} \perp \varepsilon_{2,t} \text{ no mesmo período}
\]

O IRF estrutural $\Theta_h = \Phi_h B$ mede a resposta de $y_{k,t+h}$
a um choque unitário em $\varepsilon_{j,t}$.

% ─────────────────────────────────────────────────────────────────────────────
\section{Verificação por DGP}

DGP: VAR(1) bivariado com restrição Cholesky verdadeira (y2 responde
contemporaneamente a y1, mas não vice-versa).

\begin{lstlisting}[caption={DGP estrutural e SVAR Cholesky}]
seed(42)
generate df t  = _n
// VAR(1): y1 e y2 com coeficientes cruzados
...
replace df y1 = 0.5*ly1 + 0.3*ly2 + e1t          if t == it
replace df y2 = 0.2*ly1 + 0.4*ly2 + 0.6*e1t + e2t if t == it

let m_var  = var(df, y1, y2, lags=1)
display m_var

let m_svar = svar(df, y1, y2, lags=1, id=cholesky)
display m_svar

let m_sirf = sirf(m_svar, steps=10)
display m_sirf
\end{lstlisting}

\subsection*{VAR reduzido}

\begin{lstlisting}[language={}, basicstyle=\ttfamily\small,
  frame=none, numbers=none, backgroundcolor=\color{codbg},
  xleftmargin=0pt, xrightmargin=0pt, breaklines=false]
Vector Autoregression (VAR(1)):  Observations=299  AIC=-2.83  BIC=-2.76
Residual Covariance (Σ_u):
[ 0.2452   0.1617 ]
[ 0.1617   0.3375 ]
\end{lstlisting}

\subsection*{SVAR — Matriz B}

\begin{lstlisting}[language={}, basicstyle=\ttfamily\small,
  frame=none, numbers=none, backgroundcolor=\color{codbg},
  xleftmargin=0pt, xrightmargin=0pt, breaklines=false]
============================= SVAR(1) — Cholesky =============================
B Matrix (impacto estrutural):
[ 0.4952   0.0000 ]     <- choque 1 afeta y1, não afeta y2 contemporaneamente
[ 0.3265   0.4805 ]     <- choque 2 afeta y2; y2 responde ao choque 1
\end{lstlisting}

A coluna nula em $B_{12}$ reflete a ordenação Cholesky: o choque 2 não afeta
$y_1$ contemporaneamente. $B_{21} = 0{,}33$ estima o efeito contemporâneo do
choque 1 sobre $y_2$ (DGP verdadeiro: 0{,}6 $\times$ $\sigma_{\varepsilon_1}$).

\subsection*{IRF estrutural}

\begin{lstlisting}[language={}, basicstyle=\ttfamily\small,
  frame=none, numbers=none, backgroundcolor=\color{codbg},
  xleftmargin=0pt, xrightmargin=0pt, breaklines=false]
Impulso: Choque 1 (Var0)
  h=0:  y1=+0.495  y2=+0.327   (impacto contemporâneo)
  h=1:  y1=+0.324  y2=+0.224
  h=5:  y1=+0.063  y2=+0.044   (amortecimento progressivo)
  h=9:  y1=+0.012  y2=+0.009

Impulso: Choque 2 (Var1)
  h=0:  y1= 0.000  y2=+0.481   (restrição Cholesky: y1 não responde)
  h=1:  y1=+0.115  y2=+0.170
\end{lstlisting}

A resposta de $y_1$ ao choque 2 em $h=0$ é exatamente zero — restrição de
identificação Cholesky aplicada. No horizonte $h=1$, $y_1$ começa a responder
via o canal VAR defasado.

% ─────────────────────────────────────────────────────────────────────────────
\section{Sintaxe}

\begin{lstlisting}
svar(df, y1, y2, lags=1, id=cholesky)     // SVAR com identificação Cholesky
svar(df, y1, y2, lags=1, id=longrun)      // BQ: restrições de longo prazo
sirf(m_svar, steps=10)                    // IRF estrutural
\end{lstlisting}

\begin{center}
\begin{tabular}{ll}
\toprule
\textbf{Identificação} & \textbf{Restrições} \\
\midrule
\texttt{cholesky} & triangular inferior — restrições de curto prazo recursivas \\
\texttt{longrun}  & Blanchard-Quah — restrições de longo prazo \\
\bottomrule
\end{tabular}
\end{center}

\begin{nota}
  A ordenação das variáveis no SVAR Cholesky é crucial: a variável listada
  primeiro é tratada como contemporaneamente exógena a todas as demais.
  A escolha da ordenação deve ser guiada por teoria econômica ou por
  resultados robustos a permutações da ordem.
\end{nota}
